\documentclass[11pt]{letter}
\usepackage[margin=1in]{geometry}
\usepackage{hyperref}

\signature{Liang Wang}
\address{School of Artificial Intelligence and Automation\\Huazhong University of Science and Technology\\Wuhan 430070, P.R. China\\wangliang.f@gmail.com}

\begin{document}

\begin{letter}{Editorial Office\\Results in Physics}

\opening{Dear Editor,}

We would like to submit our manuscript entitled \textbf{``Differentiable Discrete Symplectic Cosmology: Forward Simulation and Inverse Reconstruction of CMB-like Fluctuations''} for consideration in \textit{Results in Physics}.

This paper presents a differentiable computational framework based on symplectic lattice dynamics with a $1/\ln^2(t)$ adiabatic cooling law, capable of generating temperature fluctuation fields that reproduce key statistical features of the cosmic microwave background (CMB). Our main results include:

\begin{itemize}
\item Near-Gaussian statistics from symplectic evolution, validated by a 20-run ensemble and ablation against a dissipative baseline;
\item Acoustic-like oscillation peaks emerging naturally from lattice wave propagation;
\item A fully differentiable physics engine (implemented in JAX) that enables gradient-based inverse reconstruction of 3D primordial fields from real Planck SMICA data, achieving pixel correlation $r = 0.98$;
\item A full cosmic timeline simulation from Big Bang to cosmic web, revealing a V-shape thermodynamic transition driven by the cooling law.
\end{itemize}

All simulation code and Jupyter notebooks are publicly available. The manuscript has not been published elsewhere and is not under consideration by another journal.

We believe this work is well suited for \textit{Results in Physics} given its emphasis on numerical experiments, quantitative comparison with observational data, and the novel application of differentiable programming to cosmological simulation.

Thank you for your consideration.

\closing{Sincerely,}

\end{letter}
\end{document}
